\answer{ex:15:1}{(а)~$2\cdot10^5$ на нейрон, $6\cdot10^{12}$ на схему. (б)~$\approx8$~бит/с на провод, не меньше $2.5\cdot10^4$~с $\approx7$~ч.}
\answer{ex:15:2}{Шаг $2$: соседняя корреляция $0.943$, $\lambda$ от $4.18$ до $7.3\cdot10^{-4}$, отношение $\approx5.7\cdot10^3$; шаг $4$: корреляция $0.8$, отношение $\approx94$.}
\answer{ex:15:3}{(а)~$\lambda=0.988$ и $0.808$: $82$ и $4.7$ движения. (б)~$x_s^*=0.690$, $x_f^*=0.172$, сумма $0.862$.}
\answer{ex:15:4}{(а)~$116$ движений. (б)~$19$. Модель с одним процессом при нулевой сумме экономии не даёт.}
\answer{ex:15:5}{(а)~$G=50$~с$^{-1}$ против предела $10.5$~с$^{-1}$ без модели. (б)~Нет: при $G=50$~с$^{-1}$ критическая задержка $d_c\approx103\ms$; при $d=150\ms$ нужно $\hat k>0.445k$ (при любых $G$ и $d$ --- $\hat k>k/2$).}
\answer{ex:15:6}{Ошибка падает ниже $0.5\%$ за $6$--$30$~с, пол $\approx(6$--$9)\cdot10^{-4}$ при наилучших весах $7.5\cdot10^{-4}$; веса при $\eta=50$ ещё отличаются от наилучших на величину до $0.2$ вдоль плохо обусловленных направлений $C$ (задача~\ref{ex:15:2}).}
\answer{ex:15:7}{$\mathbf B=A\mathbf K$; $q_f/R\approx0.035$, $q_s/R\approx8.6\cdot10^{-4}$; при $4q_s$ $B_f\approx0.088$, $B_s\approx0.047$: медленный процесс учится более чем вдвое быстрее, а быстрый --- медленнее.}
